% Venue-specific submission source for Simulation Modelling Practice and Theory.
% The venue-neutral scientific master remains research/manuscript/manuscript.md.
\documentclass[preprint,12pt]{elsarticle}
\usepackage{amsmath,amssymb,booktabs,graphicx,url}
\journal{Simulation Modelling Practice and Theory}
\begin{document}
\begin{frontmatter}
\title{Decision-Time Validity and Cutoff Geometry in a Fixed-Budget Adaptive Digital-Twin Pipeline}
\author[aff1]{AUTHOR NAME}
\affiliation[aff1]{organization={AFFILIATION},addressline={ADDRESS},city={CITY},state={STATE},postcode={POSTCODE},country={COUNTRY}}
\begin{abstract}
Adaptive decision pipelines can exhibit strong statistical patterns that do not necessarily identify the mechanism that produced them or provide information available at decision time. We examine two such problems in a simulated adaptive digital-twin pipeline: localization of ranking changes near a fixed top-N cutoff and discrimination of harmful action-space expansion outcomes. In Experiment 166, exclusion-membership changes were strongly concentrated near the cutoff in the poisoning condition, but a prospectively frozen label-preserving non-poison perturbation matched the ranking perturbation magnitude and reproduced essentially the same localization pattern. The poison-minus-control localization difference was 0.001845 with a 95\% seed-bootstrap interval of [0.0000, 0.00554], so the frozen poisoning-specificity criterion was not met. A separate compact harmful-expansion model achieved approximately 0.979 ROC AUC retrospectively, but two of its three features incorporate post-outcome quantities and therefore cannot support a pre-decision prediction claim. Models restricted to temporally legitimate loss-surface information showed weaker exploratory discrimination. Together, the results support a pipeline-specific cutoff-localization phenomenon and motivate strict separation of observed association, perturbation specificity, causal mechanism, and prediction-time validity.
\end{abstract}
\begin{keyword}simulation validation \sep digital twin \sep decision-time validity \sep top-k selection \sep ranking stability \sep data leakage \sep falsification\end{keyword}
\end{frontmatter}
\section{Introduction}
Adaptive digital-twin research increasingly couples continuously updated digital representations with prediction, decision support, and adaptive control. Implementations span self-adaptive manufacturing, adaptive control, and digital-twin-supported decision systems \cite{splettstoesser2023selfadaptive,qiu2025digitaltwins,builesmontano2025digitaltwin}. Predictive accuracy and downstream decision quality are not interchangeable objectives. Predict-then-optimize and decision-focused learning study settings in which predicted quantities are subsequently used inside optimization or discrete decisions \cite{elmachtoub2022smart,elmachtoub2020decisiontrees,mandi2022decisionfocused}. Fixed-budget ranking creates a sensitivity question near the selection boundary: top-K theory identifies separation around the K-th and (K+1)-th items as central to reliable set identification \cite{chen2015spectral}, while ranking-stability work studies how scoring changes alter ranked outputs \cite{asudeh2018stable}. These results make near-boundary sensitivity a plausible competing explanation, not proof of the cause of Experiment 166. Leakage literature provides methodological context for unavailable-at-decision-time features \cite{kapoor2023leakage}; leakage here is established by internal feature definitions and timing audit.
\section{Methods}
\subsection{Evidence and claim governance} Analyses follow committed chronology. Preregistered criteria remain historical results; later structural and falsification analyses are subsequent adjudications. Negative findings constrain current interpretation.
\subsection{Experiment 166} Experiment 166 evaluated exclusion-membership changes in a fixed top-N pipeline using frozen features \texttt{action\_2}, \texttt{action\_3}, and \texttt{context\_support\_distance}, a class-balanced logistic-regression hazard model, 20\% targeted unsafe-to-safe source-label concealment, equal exclusion count within target seed, and 40 fresh target generation seeds (44791--44830). A switch was a candidate excluded by exactly one model. The primary near band was the closest 10\% by absolute clean cutoff margin. Criterion 1 used seed-stratified Mantel--Haenszel inference; Criterion 2 used per-seed switch-related quantities and 10,000 whole-seed bootstrap resamples. Poison mean Jaccard was 0.923823 with 308 switches. Criterion 1 yielded OR 10.567477, 95\% CI [8.345537,13.380992]; Criterion 2 yielded Spearman $\rho=-0.873179$, bootstrap 95\% CI [-0.946362,-0.735018].

Subsequent audits used a bookkeeping-preserving permutation null and a near-only versus far-only downstream-specificity analysis. A prospectively frozen stronger label-preserving control perturbed only source-training \texttt{context\_support\_distance}; 16 noise levels crossed with 16 replicates produced 256 candidates. Selection used only Jaccard and switch-count mismatch, never target outcomes or localization. The selected control had Jaccard 0.924853 and 304 switches. The frozen specificity estimand was the paired poison-minus-control difference in per-seed near-minus-far switch enrichment, with 10,000 whole-seed bootstrap resamples. Poison enrichment was 0.13623, control enrichment 0.13438, difference 0.001845, 95\% CI [0.0000,0.00554]. Strictly positive interval support was required.
\subsection{Harmful-expansion audit} The compact model used \texttt{predicted\_loss\_floor}, \texttt{loss\_floor\_error}, and \texttt{expanded\_action\_loss\_error}, where the latter two depend on \texttt{true\_best\_loss} and \texttt{realized\_expanded\_action\_loss}. The 65-event population contained 15 harmful and 50 beneficial outcomes. Seed-held-out performance was approximately 0.95 balanced accuracy, 1.00 recall, 0.75 precision, and 0.97867 AUC. The timing audit established that two residual features are post-outcome. Temporally legitimate loss-surface models showed pooled AUC approximately 0.683--0.711 and balanced accuracy 0.663--0.763; some folds contained no harmful events.
\subsection{Inferential units and provenance} Generation seed is the bootstrap unit; action rows are nested observations. The stronger-control output is preserved as a GitHub Actions artifact whose digest is verified by the manuscript plotting source before rendering.
\section{Results}
\subsection{Cutoff localization and specificity} Near-cutoff localization was strong under poisoning, but the magnitude-matched non-poison perturbation reproduced nearly the same enrichment. The paired CI did not lie strictly above zero, so poisoning specificity was not established.
\begin{figure}[t]\centering\includegraphics[width=0.92\linewidth]{generated/fig1_experiment166_poison_vs_matched_control_enrichment.png}\caption{Mean seed near-minus-far membership-switch enrichment under poisoning and the prospectively selected matched non-poison control. Rendered only from the digest-verified preserved artifact.}\label{fig:enrichment}\end{figure}
\begin{figure}[t]\centering\includegraphics[width=0.92\linewidth]{generated/fig2_experiment166_specificity_contrast.png}\caption{Frozen paired poisoning-specificity contrast with 95\% whole-seed bootstrap interval. The lower endpoint is zero; the preregistered specificity-support rule therefore was not met.}\label{fig:specificity}\end{figure}
\subsection{Mechanism adjudication} The Criterion-2 correlation was reproduced by a bookkeeping-preserving null; observed $\rho\approx-0.87318$ had one-sided null probability 0.53315. Selected-action-change rate was 0.63871 for near-only and 0.95425 for far-only contexts; difference -0.31554, 95\% CI [-0.40818,-0.23161].
\subsection{Harmful expansion} The compact model's AUC $\sim0.979$ is high retrospective discrimination, not pre-decision prediction. Temporally legitimate models remain exploratory.
\begin{table*}[t]\centering\caption{Balanced claim-adjudication summary; no new analysis.}\label{tab:adjudication}\begin{tabular}{p{0.18\textwidth}p{0.25\textwidth}p{0.31\textwidth}p{0.20\textwidth}}\toprule Claim/signal & Favorable evidence & Later discriminating evidence & Current status\\\midrule Near-cutoff localization & OR 10.567477; 95\% CI [8.345537,13.380992] & Matched non-poison control reproduced enrichment & Supported within tested pipeline\\ Poisoning specificity & First observed under poisoning & Difference 0.001845; 95\% CI [0.0000,0.00554] & Unresolved\\ Criterion-2 mechanism & $\rho=-0.873179$ & Structural null probability 0.53315 & Not independent mechanism evidence\\ Near-only downstream role & Boundary interpretation & Near-only 0.63871 vs far-only 0.95425 & Falsified in tested analysis\\ Compact prediction & AUC $\sim0.97867$ & Two post-outcome residual features & Retrospective only\\ Legitimate features & AUC $\sim0.683$--0.711 & Small/imbalanced population & Exploratory only\\\bottomrule\end{tabular}\end{table*}
\section{Negative and Failed Results} Experiment 165 failed to prospectively replicate an earlier recall $>$ AP $>$ AUC hierarchy. Experiment 158 did not satisfy preregistered prediction-decision divergence criteria. Support-distance representations performed poorly as standalone detectors, and conditioned transfer failed in at least one action/block setting. Experiment 166 mechanism and downstream-specificity interpretations failed later discriminating tests; poisoning specificity remained unresolved; and the headline harmful-expansion pre-decision interpretation contains timing leakage.
\section{Discussion} The strongest defensible Experiment 166 interpretation is pipeline-specific cutoff localization, not a poisoning-specific signature. Prior top-K and stability work makes ranking geometry a relevant competing explanation \cite{chen2015spectral,asudeh2018stable}, but does not establish mathematical necessity. The Criterion-2 association is not independent mechanistic evidence, and the downstream-specificity result runs opposite the preferential-near-switch hypothesis. High retrospective harmful-expansion discrimination coexists with invalid pre-decision timing; the non-leaking models are the relevant prospective evidence and remain exploratory. Conclusions are limited to the studied simulation, ranking rule, budget, models, perturbations, and seeds; they do not establish deployment validity or cross-domain transfer.
\section{Limitations} The matched control does not prove universal absence of poisoning-specific effects. Broader cutoff-geometry claims require varying ranking rules or budgets. The harmful-expansion population is small and imbalanced. The work is simulation-based.
\section{Conclusion} The tested fixed-budget pipeline exhibits robust near-cutoff localization, but available structural and matched-control tests do not establish poisoning specificity or an independent causal boundary mechanism. Temporally legitimate loss-surface information shows exploratory harmful-expansion discrimination, while the strongest compact result is retrospective because it incorporates post-outcome residuals. Stronger claims require new discriminating evidence.
\section*{Declaration of generative AI and AI-assisted technologies in the manuscript preparation process} During preparation, the author used OpenAI ChatGPT for research-document organization, manuscript drafting/editing, literature-search support, and adversarial claim auditing. Retained numerical results, methodological descriptions, citations, and interpretations were checked against primary literature or committed artifacts, and the author takes responsibility for the article.
\section*{Data and code availability} Code, preregistrations, result summaries, audits, source maps, and figure-generation code are maintained in the project repository. The stronger matched-control artifact provenance and SHA-256 digest are recorded in the provenance audit, and figure code verifies that digest before rendering. A stable submission snapshot will be cited by immutable Git commit/release identifier before submission.
\bibliographystyle{elsarticle-num}\bibliography{references}\end{document}
